%% kapa: exam=1 round=21 year=2010 period=2 form=A subject=appraisal_law subject_ko=부동산관계법규 seq=34 number=74
%% source: Q-Net 감정평가사 기출문제 (cst00309) artlSeq=5024181 file=extracted/1008/5024181/2010년 제21회 감정평가사 1차시험 기출문제/2교시.hwp
%% conversion: kordoc 4.24.51 -> tools/gen1.py (auto); review=proofread-gpt-sol; flags=none
\begin{question}{74}
측량․수로조사 및 지적에 관한 법령상 지적측량 손해배상책임의 보장에 관한 설명으로 옳지 않은 것은?
\begin{choices}
  \item 보증보험의 가입은 지적측량업 등록을 신청한 날로부터 10일 이내에 하여야 한다.
  \item 지적측량업자의 보증보험 가입 금액은 1억원 이상 이어야 한다.
  \item 대한지적공사의 보증보험 가입 금액은 20억원 이상 이어야 한다.
  \item 지적측량업자는 보증보험 가입 증명서류를 지적측량업을 등록한 시ㆍ도지사에게 제출하여야 한다.
  \item 보증보험에 가입한 지적측량수행자가 보증보험기간의 만료로 인하여 다시 보증보험에 가입하려는 경우에는 그 보증기간 만료일까지 다시 보증보험에 가입하여야 한다.
\end{choices}
\end{question}
\answer{1}
